\section{Related Work}

\paragraph{Speculative decoding and acceptance bottlenecks.}
Speculative decoding was introduced as a lossless way to accelerate
autoregressive generation by verifying draft tokens in parallel with the target
model \citep{leviathan2023speculative}. Subsequent work has improved the draft,
the verifier, or the scheduling policy. Online Speculative Decoding adapts the
draft to the observed query distribution to increase token acceptance
\citep{liu2023online}. Direct alignment methods train draft models to better
match chat-capable targets on plausible data distributions
\citep{goel2024direct}. AdaEDL predicts when low acceptance is likely and stops
drafting early \citep{agrawal2024adaedl}. Judge Decoding argues that model
alignment alone is insufficient and redesigns verification to improve acceptance
\citep{bachmann2025judge}. Alignment-augmented decoding uses training-free
alignment sampling and flexible verification \citep{wang2025alignment}, while
HSD pushes verification beyond token-wise acceptance by performing lossless
hierarchical verification \citep{zhou2026hsd}.

\paragraph{Why our setting is different.}
These papers all treat low acceptance as a performance obstacle to be fixed.
Our viewpoint is the inverse: low acceptance is the attack objective. Rather
than training a better draft or verifier, we ask how a user-controlled prompt
can maximize draft-target disagreement under the deployed system. To our
knowledge, prior speculative-decoding work does not explicitly optimize prompt
prefixes to minimize early acceptance probabilities.

\paragraph{Prompt-based inference-cost attacks.}
The closest attack-style prior art is prompt-based latency manipulation in
ordinary autoregressive decoding. Engorgio constructs prompts that suppress
\texttt{<EOS>} and induce abnormally long outputs, thereby increasing inference
cost and threatening service availability \citep{dong2024engorgio}. This is
related in spirit but differs in mechanism. Engorgio attacks output length in
standard autoregressive generation. Our setting attacks the verification rule in
speculative decoding by making draft proposals difficult to accept.

\paragraph{Adversarial prompt optimization and universal triggers.}
Outside speculative decoding, a large literature studies adversarial prompts,
universal triggers, and jailbreak suffixes. Universal adversarial triggers show
that short input-agnostic token sequences can induce systematic failures across
many examples \citep{wallace2019uat}. In aligned-chat settings, discrete and
hybrid search methods such as the universal suffix attack of
\citet{zou2023universal} and the hierarchical genetic search of
\citet{liu2023autodan} demonstrate that prompt optimization can reliably steer
model behavior toward harmful semantic outputs. Our problem is different in
both objective and mechanism. We do not aim to bypass safety alignment or
change the semantic content of the answer. Instead, we optimize prompts to
alter the \emph{systems behavior} of speculative verification itself: the
target output can remain correct while the verifier is forced to do
substantially more work.

\paragraph{Positioning.}
The attack studied here lies at the intersection of systems efficiency and
adversarial prompting. The speculative-decoding literature establishes that
throughput depends strongly on draft-target agreement, while the adversarial
prompt literature shows that short optimized triggers can reliably manipulate
model behavior. ADSD combines these perspectives: it turns verifier agreement
into the attack surface and develops a continuous-to-discrete optimization
framework for prompt-controlled inference-cost degradation.
